
\subsubsection{6.1 Findings}

\textbf{Finding 1: Trajectories exhibit clustering structure.} The silhouette score of 0.352 and bootstrap stability of 0.82 indicate that time series trajectories form cohesive, well-separated clusters that are robust to sampling variation.

\textbf{Finding 2: Discrete phase transitions are prevalent.} The 81\% changepoint detection rate indicates that most trajectories contain discrete phases rather than smooth continuous evolution.

\textbf{Finding 3: Two behavioral patterns dominate.} The four empirical clusters map to two dominant archetypes (slow\_burn and revival) rather than the five hypothesized. This suggests that trajectory patterns may be simpler than theoretical models predict.

\textbf{Finding 4: Growth dynamics matter more than peak timing.} Peak timing does not differentiate clusters (variance ratio 0.21), while growth ratio (4.74), changepoint count (11.08), and derivative variance (2.16) are discriminative.

\subsubsection{6.2 Limitations}

\textbf{Limitation 1: Proxy data.} The methodology was validated using astronomical lightcurve measurements rather than the originally intended dataset download statistics. The HuggingFace Hub API does not expose historical download time series. This validates the methodology but not domain-specific claims about dataset adoption.

\textbf{Limitation 2: Partial archetype recovery.} Only 2 of 5 proposed archetypes were recovered. The empirical four-cluster structure does not support the five-archetype taxonomy.

\textbf{Limitation 3: Single data source.} Analysis was conducted on one dataset from HuggingFace. Results may not generalize to other time series domains.

\textbf{Limitation 4: Baseline comparison interpretation.} The summary statistics baseline achieved higher silhouette scores than DTW clustering. While this is expected (baselines optimize for geometric separation rather than shape similarity), it complicates direct comparison of clustering quality.

\subsubsection{6.3 Implications}

The validated two-level methodology (PELT + DTW) provides a framework for trajectory analysis that separates phase detection from shape-based clustering. The finding that empirical structure maps to fewer archetypes than hypothesized suggests that trajectory characterization efforts may benefit from data-driven archetype discovery rather than a priori taxonomies.
